\section{Conclusion}
\label{sec:conclusion}

Under a fixed developmental budget, the cheapest resource is the one
already being wasted. Our diagnosis found the final third of the
Strict-Small epoch budget buying a lower training loss and no further
competence, the two curves parting company exactly where the schedule
keeps spending. That gap is the opening our method works through. Three
principles, timing, preservation, and alignment, turn the observation into
a recipe, and both a controlled study of seven interventions and the live
leaderboard bear it out. The recipe is short: train until the evaluation
saturates (Eq.~\ref{eq:wstar}), then spend what remains on anchored,
evaluation-aligned refinement from that checkpoint, held close enough that
the hard-won structure survives. Nothing in it is specific to our
backbone, and both steps should carry to any model whose budget outlives
its learning. We resist a strong developmental reading, though the shape
is suggestive: exposure first, contrast after. At this scale, when the
corpus is exhausted, the budget need not be.
